\documentclass{article}

\usepackage{arxiv}
\rhead{\scshape Preprint}

\usepackage[utf8]{inputenc}
\usepackage[T1]{fontenc}
\usepackage{hyperref}
\usepackage{url}
\usepackage{booktabs}
\usepackage{amsfonts}
\usepackage{amsmath,amssymb,amsthm}
\usepackage{mathtools}
\usepackage{nicefrac}
\usepackage[expansion=false]{microtype}
\usepackage{algorithm}
\usepackage{algpseudocode}
\usepackage{natbib}
\usepackage{doi}

%% ── theorem environments ──────────────────────────────────────────────────
\newtheorem{theorem}{Theorem}
\newtheorem{lemma}[theorem]{Lemma}
\newtheorem{corollary}[theorem]{Corollary}
\newtheorem{definition}{Definition}
\newtheorem{remark}{Remark}
\newtheorem{invariant}{Invariant}

%% ── macros ────────────────────────────────────────────────────────────────
\newcommand{\Ls}{\mathcal{L}^{*}}
\newcommand{\Li}[1]{\mathcal{L}^{(#1)}}   % L^{(i)}: iterate i
\newcommand{\Hi}[1]{H^{(#1)}}             % H^{(i)}: H at iterate i
\newcommand{\PL}{\mathrm{PL}}
\newcommand{\Shad}{\mathrm{Shad}}
\newcommand{\Filter}{\mathrm{Filter}}
\newcommand{\bPhi}{\boldsymbol{\Phi}}
\newcommand{\Zm}{\mathbb{Z}_m}
\newcommand{\Hstar}{H^{*}}
\newcommand{\Estar}{E^{*}}
\newcommand{\own}{\mathit{own}}
\newcommand{\writ}{\mathit{written}}
\newcommand{\pred}{\mathit{pred}}
\DeclareMathOperator{\ord}{ord}

\newcommand{\Lstar}{L^{*}}
\newcommand{\N}{\mathbb{N}}

%% ── title block ───────────────────────────────────────────────────────────
\title{\textbf{Decidability of Livelock Detection\\
       for Parameterized Self-Disabling Unidirectional Rings}}

\author{
  Aly Farahat\\
  \texttt{aly.farahat@intusurg.com}
}

\begin{document}
\maketitle

\begin{abstract}
We prove that livelock detection is \emph{decidable in polynomial time}
for parameterized symmetric unidirectional rings of self-disabling
processes with bounded domain $\mathbb{Z}_m$.
Given a protocol $T$, the algorithm decides whether a livelock exists
for \emph{some} ring size $K\!\geq\!2$, running in $O(|T|^3)$ time
independent of $K$.
The algorithm computes $\Ls$---the greatest fixed point of a
deflationary monotone operator on the finite set $T$---and returns
\textsc{livelock} iff $\Ls\!\neq\!\emptyset$.
The livelock freedom argument rests on maximality: $\Ls$ is the
largest set of transitions that can simultaneously sustain a
pseudolivelock at every process; its emptiness certifies freedom for
all $K$ without any search over ring sizes.
The work is grounded in the algebraic characterization of livelocks
from Farahat~\citep{farahat2012}, which establishes necessary and
sufficient conditions for livelock existence but does not address
decidability. We also handle the $(1,1)$-asymmetric case in which one
distinguished process $P_0$ differs from the remaining $K\!-\!1$
identical processes.
Code and algebraic foundation are at
\url{https://github.com/cosmoparadox/mathematical-tools}.
\end{abstract}

\keywords{livelock detection \and self-stabilization \and
          parameterized verification \and unidirectional ring \and
          fixed point}

%% ═══════════════════════════════════════════════════════════════════════════
\section{Introduction}
%% ═══════════════════════════════════════════════════════════════════════════

Self-stabilizing distributed protocols recover from arbitrary transient
faults by converging to a legitimate state from any initial
configuration~\citep{dijkstra1974}. A key correctness property is
\emph{livelock freedom}: no infinite execution in which every process
fires indefinitely without stabilizing. For protocols on a unidirectional
ring of $K$ processes, the verification question is whether livelock
freedom holds for all $K\!\geq\!2$---the \emph{parameterized} livelock
freedom problem.

Klinkhamer and Ebnenasir~\citep{klinkhamer2013,klinkhamer2019} study
this problem for self-disabling deterministic protocols, establishing
$\Pi^0_1$-completeness and a correct semi-algorithm for livelock
\emph{detection}.  Building on the algebraic characterization of
Farahat~\citep{farahat2012}, where pseudolivelock and propagation are
structured computable relations on $T$, the present paper shows that
livelock \emph{freedom} is decidable in $O(|T|^3)$ time via the
greatest fixed point $\Ls$ of a single operator on $T$, independent of
ring size $K$.

\paragraph{Contribution.}
We show livelock freedom reduces to checking $\Ls = \emptyset$, where
$\Ls$ is the greatest fixed point of a single deflationary operator on
$T$. This runs in $O(|T|^3)$ time independent of $K$. It suffices that
\emph{one} process has no pseudolivelock to certify freedom for all
ring sizes simultaneously.

\paragraph{Foundation.}
This paper builds on Farahat~\citep{farahat2012} (Ch.~6), which
characterizes livelocks via the \emph{Propagation Law}
$H\!\circ\!E = E\!\circ\!H$ and the \emph{Circulation Law}
$H^{|E|}\cap E^K\!\neq\!\emptyset$. That work is existential and does
not address decidability. The present paper makes it computationally
effective. Code and the dissertation chapter are at
\url{https://github.com/cosmoparadox/mathematical-tools}.

%% ═══════════════════════════════════════════════════════════════════════════
\section{Setting and Definitions}
%% ═══════════════════════════════════════════════════════════════════════════

\paragraph{Ring topology.}
A \emph{unidirectional ring} of size $K\!\geq\!2$ consists of processes
$P_0,P_1,\ldots,P_{K-1}$ with indices mod $K$. Each process $P_r$ has a
state variable $x_r \in \Zm = \{0,1,\ldots,m-1\}$ for a fixed integer
$m\!\geq\!2$. Process $P_r$ reads its \emph{predecessor value}
$x_{r-1}$ and its \emph{own value} $x_r$, and may write a new value to
$x_r$.

\begin{definition}[Transition and protocol]
\label{def:trans}
A \emph{transition} is a triple $t=(v,w,w')\in\Zm^3$ where:
\begin{itemize}
  \item $v = t.\pred$: the predecessor value that enables $t$,
  \item $w = t.\own$: the own value that enables $t$,
  \item $w' = t.\writ$: the value written to $x_r$ when $t$ fires.
\end{itemize}
A \emph{protocol} is a finite set $T\subseteq\Zm^3$ of transitions.
A transition $t=(v,w,w')\in T$ satisfies $w'\neq w$ (the written value
differs from the own value). It is \emph{self-disabling} if no
transition in $T$ has the form $(v, w', w'')$ for any $w''\in\Zm$;
that is, no transition in $T$ shares the same predecessor value $v$
and has own value equal to the written value $w'$ of $t$. Intuitively,
after $t$ fires and sets $x_r\!\leftarrow\!w'$, process $P_r$ cannot
fire again even if the predecessor value stays $v$, because no
transition is enabled with $(\pred,\own) = (v,w')$.
Throughout this paper, \emph{all transitions in $T$ are assumed
self-disabling}.
A \emph{protocol} is \emph{self-disabling} if every transition in it
is self-disabling.
\end{definition}

\begin{definition}[Enabling and execution]
\label{def:exec}
Transition $t=(v,w,w')$ is \emph{enabled} at process $P_r$ in global
state $\mathbf{x}=(x_0,\ldots,x_{K-1})$ iff $x_{r-1}=v$ and $x_r=w$.
Firing $t$ sets $x_r\leftarrow w'$ and leaves all other variables
unchanged. An \emph{execution} is an infinite sequence of
(process, transition) pairs, each pair enabled at the time of firing.
A \emph{livelock} is an execution in which every process fires
infinitely often.
\end{definition}

\begin{definition}[$(l,q)$-ring topology]
\label{def:lq}
A \emph{$(l,q)$-ring} of size $K = l + nq$ (for integer $n\geq 1$)
is a unidirectional ring in which:
\begin{itemize}
  \item Processes $P_0,\ldots,P_{l-1}$ form a \emph{prefix} of $l$
        distinct process types, each appearing exactly once.
  \item Processes $P_l,\ldots,P_{l+q-1}$ form a \emph{period} of $q$
        distinct process types, repeated $n$ times to fill the ring.
\end{itemize}
A ring is \emph{symmetric} ($l=0$, $q=1$) when all $K$ processes run
the same protocol $T$; valid ring sizes are all $K\geq 2$.
A ring is \emph{$(1,1)$-asymmetric} ($l=1$, $q=1$) when $P_0$ runs
protocol $T_0$ and $P_1,\ldots,P_{K-1}$ all run $T_{\mathit{other}}$;
valid ring sizes are $K = 1+n$ for any $n\geq 1$.
This paper addresses the symmetric and $(1,1)$-asymmetric cases.
The general $(l,q)$-asymmetric case is left to a companion paper.
\end{definition}

\begin{definition}[Pseudolivelock graph $H(S)$ and operator $\PL$]
\label{def:H}
For a set of transitions $S\subseteq T$, define the directed graph
$H(S)$ as follows:
\begin{itemize}
  \item \emph{Vertices}: the transitions in $S$.
  \item \emph{Edges}: there is a directed edge $t_i\!\to\!t_j$ in
        $H(S)$ iff $t_j.\own = t_i.\writ$, i.e., the value written by
        $t_i$ equals the own value required to enable $t_j$.
\end{itemize}
The operator $\PL(S)$ returns the set of transitions in $S$ that lie on
at least one directed cycle in $H(S)$:
\[
  \PL(S) \;=\; \bigcup\, \{C \subseteq S : C \text{ is a non-trivial
  strongly connected component of } H(S)\}.
\]
Algorithmically, $\PL(S)$ is computed by Tarjan's SCC algorithm on
$H(S)$, retaining only SCCs of size $\geq 2$ and self-loops (edges
$t\to t$). Complexity: $O(|S|^2)$.
\end{definition}

\begin{definition}[Shadow of a pseudolivelock --- $\Shad$]
\label{def:shad}
Let $P\subseteq T$ be a set of transitions with pseudolivelock graph
$H(P)$.  For each directed $H(P)$-edge $t_i\!\to\!t_j$ (where
$t_j.\own = t_i.\writ$), define the \emph{arc} of that edge as the
pair $(t_i.\pred,\, t_j.\pred)$.  Intuitively: when process $P_r$
fires $t_i$ and then $t_j$ in succession, its predecessor value
$x_{r-1}$ must change from $t_i.\pred$ to $t_j.\pred$ between the
two firings.  The transition in $P_{r-1}$ responsible for this change
must have $\own = t_i.\pred$ and $\writ = t_j.\pred$.  We call
$(t_i.\pred,\, t_j.\pred)$ the \emph{shadow} of the edge
$t_i\!\to\!t_j$.

The \emph{shadow set} of $P$ collects all such required $(\own,\writ)$
pairs over every $H(P)$-edge:
\[
  \Shad(P) \;=\;
  \bigl\{\,(t_i.\pred,\; t_j.\pred) :
     t_i,t_j\in P,\; t_j.\own = t_i.\writ\,\bigr\}.
\]
A predecessor transition $t_k\in T_{r-1}$ is a \emph{shadow witness}
for the edge $t_i\!\to\!t_j$ iff $(t_k.\own,\,t_k.\writ)=(t_i.\pred,\,t_j.\pred)$,
i.e., $t_k$ has the form $(\,\cdot\,,\, t_i.\pred,\, t_j.\pred\,)$
in $T_{r-1}$.  The shadow set $\Shad(P)$ is exactly the set of
$(\own,\writ)$ pairs that must be present among the predecessor
transitions to sustain equivariant backward propagation from $P_r$ to
$P_{r-1}$.
\end{definition}

\begin{definition}[Filter operator $\Filter$]
\label{def:filter}
For a set $S\subseteq T$ and a requirement set $R\subseteq\Zm^2$:
\[
  \Filter(S,\,R) \;=\; \{\,t\in S : (t.\own,\; t.\writ)\in R\,\}.
\]
$\Filter(T_{r-1},\,\Shad(P))$ is the \emph{backward propagation step}:
it extracts from $T_{r-1}$ every transition of the form
$(\,\cdot\,,\, t_i.\pred,\, t_j.\pred\,)$ for some $H(P)$-edge
$t_i\!\to\!t_j$.  These are precisely the shadow witnesses in $T_{r-1}$
that can sustain equivariant propagation backward from $P_r$ to
$P_{r-1}$.  Any transition in $T_{r-1}$ that is not a shadow witness
for any edge in $P$ is removed, since it cannot participate in a
consistent backward propagation.
\end{definition}

\begin{remark}[Shadow witnesses are exactly the equivariance condition]
\label{rem:equiv}
We explain why the existence of shadow witnesses in $T_{r-1}$ is
equivalent to the equivariance condition
$H_{r-1}\circ E \;=\; E\circ H_r$.

Fix an $H_r$-edge $(t_i\!\to\!t_j)$ in $L_r$ and its shadow witness
$t_k\in L_{r-1}$ with $t_k.\own=t_i.\pred$ and $t_k.\writ=t_j.\pred$.
The arc construction sets $E[t_k][t_j]=1$.  Now expand both sides:

\emph{Left side $(H_{r-1}\circ E)[t_k][t_j]=1$}
requires $\exists\,t_s$: $H_{r-1}[t_k][t_s]=1$ and $E[t_s][t_j]=1$.
The first condition gives $t_s.\own = t_k.\writ = t_j.\pred$; the
second gives $t_s.\own = t_{i'}.\pred$ and $t_s.\writ = t_j.\pred$
for some $H_r$-edge $(t_{i'}\!\to\!t_j)$.  These are consistent: $t_s$
is the $H_{r-1}$-successor of $t_k$ with $t_s.\own = t_j.\pred$.

\emph{Right side $(E\circ H_r)[t_k][t_j]=1$}
requires $\exists\,t_s$: $E[t_k][t_s]=1$ and $H_r[t_s][t_j]=1$.
The first gives a $H_r$-edge $(t_i\!\to\!t_s)$ with $t_k.\own=t_i.\pred$
and $t_k.\writ=t_s.\pred$; the second gives $t_j.\own=t_s.\writ$.

\emph{Why both sides agree:}
The shared value is $t_j.\pred = t_k.\writ = t_s.\own$.  It links the
$H_{r-1}$-successor of $t_k$ (whose $\own$ must equal $t_j.\pred$) to
the $H_r$-successor of $t_i$ (whose $\pred$ must equal $t_k.\writ =
t_j.\pred$).  The shadow condition $(t_k.\own,\,t_k.\writ) =
(t_i.\pred,\,t_j.\pred)$ is therefore precisely the statement that the
arc $t_k\!\to\!t_s$ in $H_{r-1}$ mirrors the arc $t_i\!\to\!t_j$ in
$H_r$---which is exactly what $H_{r-1}\circ E = E\circ H_r$ asserts.

Consequently, $\Filter(T_{r-1},\Shad(P_r))$ retains precisely those
transitions of $T_{r-1}$ that can participate in an equivariant
propagation from $P_r$ to $P_{r-1}$.
\end{remark}

\begin{definition}[Fixed-point operator $\bPhi$ and livelock kernel $\Ls$]
\label{def:phi}
\emph{Design rationale.}
Intuitively, $\bPhi$ removes transitions that cannot participate in any
equivariant backward propagation and retains only those forming maximal
pseudolivelocks.  By Remark~\ref{rem:equiv}, one iteration implements a
backward equivariant propagation step: it (i)~extracts the maximal
pseudolivelock $\PL(S)$, (ii)~computes the shadow $\Shad(\PL(S))$
encoding the equivariance requirement at the predecessor interface, and
(iii)~filters $S$ and re-extracts the maximal pseudolivelock.

Formally, define $\bPhi : 2^T \to 2^T$ by:
\[
  \bPhi(S)
  \;=\; \underbrace{\PL\!\bigl(}_{\text{(iii) re-extract SCC}}
         \underbrace{\Filter(S,\;}_{\text{(ii) prune non-witnesses}}
         \underbrace{\Shad(\PL(S))}_{\text{(i) required pairs}}
         \bigr).
\]
Starting from $\Li{0}=T$, iterate $\Li{i+1}=\bPhi(\Li{i})$.
The \emph{livelock kernel} is the limit:
\[
  \Ls \;=\; \lim_{i\to\infty}\Li{i} \;=\; \nu S.\,\bPhi(S),
\]
the greatest fixed point of $\bPhi$.  We write $\Hi{i}=H(\Li{i})$.

\noindent\textbf{Properties used throughout:}
\begin{enumerate}
  \item[\textbf{(D)}] \emph{Deflationary:} $\bPhi(S)\subseteq S$,
        because both $\Filter$ and $\PL$ only remove transitions.
  \item[\textbf{(M)}] \emph{Monotone:} $S'\!\subseteq\!S \Rightarrow
        \bPhi(S')\!\subseteq\!\bPhi(S)$, because smaller $S$ yields
        smaller $\PL(S)$, hence smaller $\Shad(\PL(S))$, hence fewer
        transitions survive the filter.
\end{enumerate}
\end{definition}

\begin{definition}[Propagation $\Estar$]
\label{def:E}
Given $\Ls$ and the graph $\Hstar = H(\Ls)$, define the
\emph{propagation relation} $\Estar$ as follows. For each directed edge
$(t_i\to t_j)$ in $\Hstar$, find each transition $t_k\in\Ls$
satisfying:
\[
  t_k.\own = t_i.\pred \quad\text{and}\quad t_k.\writ = t_j.\pred.
\]
Set $\Estar[t_k][t_j]=1$ for each $k$ such that $t_k \in \Ls$. Each transition
$t_k$ is a \emph{predecessor enabler}: it is the transition that the
predecessor process $P_{r-1}$ must fire to make the edge $(t_i\to t_j)$
propagate forward to $P_r$.
\end{definition}

%% ═══════════════════════════════════════════════════════════════════════════
\section{Algorithm}
%% ═══════════════════════════════════════════════════════════════════════════

\begin{algorithm}[h]
\caption{Algorithm~1: Symmetric livelock decision}
\label{alg:main}
\begin{algorithmic}[1]
\Require Symmetric protocol $T$ (all $K$ processes identical)
\Ensure \textsc{livelock} or \textsc{free}
\State $S \leftarrow T$
\Repeat
  \State $P \leftarrow \PL(S)$
    \hfill\Comment{maximal SCC of $H(S)$; Def.~\ref{def:H}}
  \If{$P = \emptyset$} \Return \textsc{free} \EndIf
  \State $S' \leftarrow \PL(\Filter(S,\,\Shad(P)))$
    \hfill\Comment{shadow back-propagation; Defs.~\ref{def:shad},\ref{def:filter}}
\Until{$S' = S$}
\State $\Ls \leftarrow S'$
\Return \textsc{livelock}
\end{algorithmic}
\end{algorithm}

\begin{algorithm}[h]
\caption{Algorithm~2: $(1,1)$-Asymmetric livelock decision}
\label{alg:asym}
\begin{algorithmic}[1]
\Require Protocols $T_0$ ($P_0$) and $T_{\mathit{other}}$ ($P_1,\ldots,P_{K-1}$)
\Ensure \textsc{livelock} or \textsc{free}
\State $L_0 \leftarrow \PL(T_0)$;\quad
       $L_{\mathit{other}} \leftarrow \PL(T_{\mathit{other}})$
\If{$L_0=\emptyset$ \textbf{or} $L_{\mathit{other}}=\emptyset$}
  \Return \textsc{free}
\EndIf
\While{true} \hfill\Comment{joint fixed-point over $(L_0,L_{\mathit{other}})$}
  \State \Comment{Step 1: constrain $L_{\mathit{other}}$ using current $L_0$,
         then re-stabilise with Algorithm~1}
  \State $L_{\mathit{other}} \leftarrow
    \text{Algorithm~1 warm-started from }
    \PL(\Filter(L_{\mathit{other}},\Shad(L_0)))$
  \If{$L_{\mathit{other}}=\emptyset$} \Return \textsc{free} \EndIf
  \State \Comment{Step 2: ring closure back to $P_0$ using updated $L_{\mathit{other}}$}
  \State $L_0' \leftarrow \PL(\Filter(T_0,\,\Shad(L_{\mathit{other}})))$
  \If{$L_0'=\emptyset$} \Return \textsc{free} \EndIf
  \If{$L_0' = L_0$} \textbf{break}
    \hfill\Comment{joint fixed point reached}
  \EndIf
  \State $L_0 \leftarrow L_0'$
    \hfill\Comment{update $L_0$ before next iteration}
\EndWhile
\State $\Ls \leftarrow L_0'$;\quad $\Ls_{\mathit{other}} \leftarrow L_{\mathit{other}}$
\Return \textsc{livelock}
\end{algorithmic}
\end{algorithm}

\paragraph{Notation for the proof.}
Let $\Li{0}=T$ and $\Li{i+1}=\bPhi(\Li{i})$ be the iterates produced
by Algorithm~\ref{alg:main}. Write $\Hi{i}=H(\Li{i})$ for the
pseudolivelock graph at iterate $i$.

%% ═══════════════════════════════════════════════════════════════════════════
\section{Correctness}
\label{sec:correct}

The proof uses two results from Farahat as black boxes.

\begin{lemma}[Farahat's characterisation~{\citep[Ch.~6]{farahat2012}}]
\label{lem:farahat}
A livelock exists for some $K\!\geq\!2$ if and only if there exist a
pseudolivelock $H$ and propagation $E$ on a non-empty support
$\mathcal{L}\subseteq T$, satisfying:
\begin{enumerate}
  \item[\textbf{(F1)}] Every $t\in\mathcal{L}$ lies on a directed cycle
        in $H(\mathcal{L})$.
  \item[\textbf{(F2)}] $E$ is defined on all of $\mathcal{L}$: every
        $t_k\in\mathcal{L}$ is the arc predecessor of some $H$-edge
        $(t_i\!\to\!t_j)$ in $\mathcal{L}$, i.e., participates in at
        least one propagation arc. Here, $E$ is strengthened by the equivariance condition with the predecessor's interface.
  \item[\textbf{(F3)}] $H\circ E = E\circ H$ at every process
        interface.
\end{enumerate}
We call such $(H,E,\mathcal{L})$ a \emph{Farahat witness}.
\end{lemma}

\begin{lemma}[Circulation Law~{\citep[Ch.~6]{farahat2012}}]
\label{lem:cl}
Any Farahat witness yields a concrete ring size $K'\geq 2$ and a
periodic livelock execution on a ring of size $K'$.
\end{lemma}

The theorem has two directions with different proof strategies.
\emph{Soundness} ($\Ls\neq\emptyset\Rightarrow$ livelock): the algorithm
returns a non-empty fixed point $\Ls$ on which it constructs a
pseudolivelock $\Hstar$ at every process and a propagation $\Estar$
satisfying equivariance; Farahat's theorem then gives the livelock.
This direction is straightforward: we simply verify (F1)--(F3) for the
objects $\Hstar,\Estar,\Ls$ built by the fixed point.
\emph{Completeness} (livelock $\Rightarrow\Ls\neq\emptyset$): if a
livelock exists, Farahat guarantees a witness $(H,E,\mathcal{L})$; we
show $\mathcal{L}$ is a pre-fixed point of $\bPhi$, so by maximality of
$\Ls$ (Knaster-Tarski), $\mathcal{L}\subseteq\Ls$ and $\Ls\neq\emptyset$.
This direction is harder: it requires showing that the equivariance
conditions enforced by $\bPhi$ are exactly the ones satisfied by every
valid witness, so no witness support can escape the fixed point.

The following invariant is central to both directions.

\begin{invariant}[Maximality of $\Hi{i}$]
\label{inv:max}
At every iterate $i\geq 0$, every transition in $\Li{i}$ lies on a
directed cycle in $\Hi{i}=H(\Li{i})$.
\end{invariant}

\begin{proof}[Proof of Invariant~\ref{inv:max}]
By induction on $i$.
\emph{Base ($i=0$):} $\Li{1}=\bPhi(T)$.
Every element of $\Li{1}$ is in $\PL(\cdot)$ by definition, so lies on
a cycle in $H(\Li{1})=\Hi{1}$.
\emph{Step:} Assume the invariant at $i$, so $\PL(\Li{i})=\Li{i}$.
Then $\Li{i+1}=\PL(\Filter(\Li{i},\Shad(\Li{i})))$; every element of
$\Li{i+1}$ lies on a cycle in $\Hi{i+1}$ by definition of $\PL$. \qed
\end{proof}

\begin{lemma}[Completeness: every Farahat witness support is a pre-fixed point of $\bPhi$]
\label{lem:prefixed}
Let $(H,E,\mathcal{L})$ be a Farahat witness.  Then
$\mathcal{L}\subseteq\bPhi(\mathcal{L})$; hence $\mathcal{L}\subseteq\Ls$.
\end{lemma}

\begin{proof}
We verify $\mathcal{L}\subseteq\PL(\mathcal{L})$ and
$\mathcal{L}\subseteq\Filter(\mathcal{L},\Shad(\PL(\mathcal{L})))$.

\emph{$\mathcal{L}\subseteq\PL(\mathcal{L})$:}
By hypothesis (F1) of Lemma~\ref{lem:farahat}, $H$ is a pseudolivelock
on $\mathcal{L}$: every $t\in\mathcal{L}$ lies on a directed cycle in
$H=H(\mathcal{L})$.  Hence $\mathcal{L}\subseteq\PL(\mathcal{L})$, and
in fact $\PL(\mathcal{L})=\mathcal{L}$.

\emph{$\mathcal{L}\subseteq\Filter(\mathcal{L},\Shad(\mathcal{L}))$:}
By hypothesis (F2), $E$ maps $\mathcal{L}$ into $\mathcal{L}$: for each
$H$-edge $(t_i\!\to\!t_j)$ in $\mathcal{L}$, the arc predecessor
$t_k\in\mathcal{L}$ with $t_k.\own=t_i.\pred$ and $t_k.\writ=t_j.\pred$
exists.  Therefore $(t_k.\own,t_k.\writ)=(t_i.\pred,t_j.\pred)\in\Shad(\mathcal{L})$,
so $t_k\in\Filter(\mathcal{L},\Shad(\mathcal{L}))$.  If $t_k$
did not participate in any propagation arc, then $E$ would not be
defined at $t_k$, contradicting hypothesis~(F2) that $(H,E,\mathcal{L})$
is a valid Farahat witness.  Hence every $t_k\in\mathcal{L}$ is the arc
predecessor of some $H$-edge, satisfying
$(t_k.\own,t_k.\writ)\in\Shad(\mathcal{L})$ and giving
$\mathcal{L}\subseteq\Filter(\mathcal{L},\Shad(\mathcal{L}))$.
(The operator $\bPhi$ enforces this condition by construction, pruning any
transition whose $(\own,\writ)$ pair does not appear as a shadow
witness.)\looseness=-1

Combining: $\mathcal{L}\subseteq\PL(\Filter(\mathcal{L},\Shad(\mathcal{L})))
=\bPhi(\mathcal{L})$.  Since $\bPhi$ is monotone and $\Ls$ is the
greatest fixed point of $\bPhi$, every pre-fixed point is contained in
$\Ls$ (Knaster-Tarski~\citep{tarski1955}); hence $\mathcal{L}\subseteq\Ls$. \qed
\end{proof}

\begin{lemma}[Soundness: $\Ls\neq\emptyset$ implies the algorithm constructs a Farahat witness]
\label{lem:witness}
If $\Ls\neq\emptyset$, then $(\Hstar,\Estar,\Ls)$ is a Farahat witness.
\end{lemma}

\begin{proof}
We verify hypotheses (F1)--(F3) of Lemma~\ref{lem:farahat}.

\textbf{(F1)}~Invariant~\ref{inv:max} is proved by induction on the
iteration index $i$: at every step, $\bPhi$ applies $\PL$, which by
definition retains only transitions lying on directed cycles, so
$\Hi{i}$ is maximal at every iterate.  At the fixed point
$\Ls=\Li{i^*}$, the invariant gives: every $t\in\Ls$ lies on a
directed cycle in $\Hstar=H(\Ls)$, establishing~(F1).

\textbf{(F2)}~The fixed-point equation $\bPhi(\Ls)=\Ls$ gives
$\Ls=\PL(\Filter(\Ls,\Shad(\Ls)))$, so $\Shad(\Ls)$ is satisfied within
$\Ls$: for each $\Hstar$-edge $(t_i\!\to\!t_j)$, the pair
$(t_i.\pred,t_j.\pred)\in\Shad(\Ls)$, guaranteeing at least one
$t_k\in\Ls$ with $t_k.\own=t_i.\pred$ and $t_k.\writ=t_j.\pred$.  Hence $\Estar$ is
defined on all of $\Ls$ and every transition participates in a
propagation arc.

\textbf{(F3)}~By Remark~\ref{rem:equiv}, the shadow-witness condition
of (F2) is algebraically equivalent to
$\Hstar\circ\Estar=\Estar\circ\Hstar$. \qed
\end{proof}

\begin{lemma}[Termination and maximality]
\label{lem:termination}
Algorithm~\ref{alg:main} terminates in at most $|T|$ iterations in
$O(|T|^3)$ time.  $\Ls$ is the greatest subset of $T$ closed under
shadow back-propagation and pseudolivelock completion; it therefore
contains every support admitting a Farahat witness.
\end{lemma}

\begin{proof}
$\bPhi$ is deflationary~(D) and monotone~(M), so
Knaster-Tarski~\citep{tarski1955} gives a greatest fixed point in at
most $|T|$ steps (each removes $\geq 1$ transition).  Each step:
$O(|T|^2)$ Tarjan $+$ $O(|T|)$ filter; total $O(|T|^3)$, independent
of $K$.  By Lemma~\ref{lem:prefixed}, every Farahat witness support
$\mathcal{L}$ satisfies $\mathcal{L}\subseteq\bPhi(\mathcal{L})$,
making it a pre-fixed point of $\bPhi$; by Knaster-Tarski,
$\mathcal{L}\subseteq\Ls$. \qed
\end{proof}

\begin{theorem}[Main result]
\label{thm:main}
Algorithm~\ref{alg:main} decides in $O(|T|^3)$ time, independent of $K$:
\[
  \Ls \neq \emptyset
  \;\iff\;
  \text{a livelock exists for some } K \geq 2.
\]
\end{theorem}

\begin{proof}
\emph{Soundness} $(\Rightarrow)$: Lemma~\ref{lem:witness} constructs the
Farahat witness $(\Hstar,\Estar,\Ls)$; Lemma~\ref{lem:cl} gives the
livelock.
\emph{Completeness} $(\Leftarrow)$: any livelock yields a Farahat witness
$(H,E,\mathcal{L})$ by Lemma~\ref{lem:farahat}; Lemma~\ref{lem:prefixed}
gives $\mathcal{L}\subseteq\Ls$, so $\Ls\neq\emptyset$.
Complexity: Lemma~\ref{lem:termination}. \qed
\end{proof}

\paragraph{Livelock freedom---the maximality argument.}
When $\Ls=\emptyset$, Algorithm~\ref{alg:main} certifies livelock
freedom for \emph{all} $K\!\geq\!2$ in one $O(|T|^3)$ computation.
By Lemma~\ref{lem:termination}, $\Ls$ contains every Farahat witness support.
Its emptiness means no subset of $T$ is a pre-fixed point of $\bPhi$:
there is no set of transitions that can simultaneously sustain a
pseudolivelock at every process with consistent predecessor propagation.
Crucially, it suffices that \emph{one} process has an empty
pseudolivelock---because Invariant~\ref{inv:max} ensures that at every
iterate $i$, every transition in $\Li{i}$ lies on a cycle in $\Hi{i}$.
If at any iterate $P=\PL(\Li{i})=\emptyset$, the algorithm terminates
immediately, having proved that no ring of any size $K$ can sustain a
livelock.

%======================================================
\section{Semantic Faithfulness of the Livelock Kernel}
\label{sec:semantics}
Theorem~\ref{thm:main} established $\Ls\neq\emptyset \iff$
livelock via Farahat's characterisation (Lemma~\ref{lem:farahat})
as a black box: the proof never reasons about executions directly.
This section provides a structural explanation of why livelock
witnesses are bounded and does not rely on the fixed-point algorithm.
The argument proceeds entirely through the execution semantics
of livelocks, yielding an independent proof of decidability
(Corollary~\ref{cor:decidability}) and an independent proof that
$\Ls$ is semantically faithful
(Theorem~\ref{thm:faithful}).

The key structural insight is a hierarchy of reduction:
\begin{enumerate}
  \item[\textbf{Level~1.}] \emph{General:} A livelock induces
    closed equivariant walks at every process
    (Remark~\ref{rem:periodic}).
  \item[\textbf{Level~2.}] \emph{Reduced:} Every such walk
    decomposes into simple equivariant cycles of bounded length
    (Lemmas~\ref{lem:decompose}--\ref{lem:walk-simple}).
  \item[\textbf{Level~3.}] \emph{Structural:} These simple cycles
    propagate bijectively (up to cyclic rotation) across processes,
    and the Circulation Law reduces to finite arithmetic over
    cycle propagation and shift arithmetic
    (Section~\ref{sec:circulation}).
\end{enumerate}

\subsection{Livelock executions and periodic walks}

\begin{remark}[Periodic structure of livelocks {\citep{klinkhamer2019}}]\label{rem:periodic}
If a livelock exists on a ring of size~$K$, it admits
a periodic execution in which each process repeatedly
traverses a closed walk in~$\Hstar(\Ls)$.
These walks are \emph{equivariant}: the enabling walk for
each process's closed walk is the closed walk of its
predecessor.
\end{remark}

\subsection{Completeness: livelock implies
\texorpdfstring{$\Ls \neq \emptyset$}{L* nonempty}}

We show that any livelock execution produces a nonempty
fixed point of~$\bPhi$, and that the witness consists of
simple cycles.

\begin{lemma}[Closed walks contain simple cycles]
\label{lem:decompose}
Every closed walk in a directed graph contains a simple cycle.
\end{lemma}

\begin{proof}
If any transition appears more than once in the walk, the
sub-walk between the two occurrences is a shorter closed walk.
Iterating yields a closed walk with no repeated transitions,
which is a simple cycle.
\end{proof}

\begin{lemma}[Simple cycles induce equal-length closed walks]
\label{lem:equal-length}
Let $C$ be a simple cycle of length~$N$ in~$\Hstar(\Ls)$.
Every enabling walk for~$C$ is a closed walk of length~$N$.
\end{lemma}

\begin{proof}
Each of the~$N$ edges of~$C$ contributes exactly one shadow
constraint, requiring exactly one witness transition.
The enabling walk therefore has exactly~$N$ entries.
Closure follows from the cyclicity of~$C$:
the shadow of the last edge $(t_{N-1}, t_0)$ requires
$w_{N-1}.\mathit{written} = t_0.\mathit{pred}$,
while the shadow of the first edge $(t_0, t_1)$ requires
$w_0.\mathit{own} = t_0.\mathit{pred}$.
Hence $w_{N-1}.\mathit{written} = w_0.\mathit{own}$,
which is the H-edge condition $w_{N-1} \to w_0$,
closing the walk.
\end{proof}

\begin{definition}[Shadow obligation chain]
\label{def:shadow-chain}
Let $C = (t_0, t_1, \ldots, t_{N-1})$ be a simple cycle of length~$N$
in~$\Hstar(\Ls)$.
The \emph{shadow obligation chain} of~$C$ is the cyclic sequence of
shadow pairs:
\[
  (t_0.\pred,\, t_1.\pred),\;
  (t_1.\pred,\, t_2.\pred),\;
  \ldots,\;
  (t_{N-1}.\pred,\, t_0.\pred).
\]
Each pair constrains the predecessor: position~$k$ of the enabling walk
must satisfy $w_k.\own = t_k.\pred$ and $w_k.\writ = t_{k+1}.\pred$.
The chain is \emph{irreducible} if no proper contiguous sub-sequence
closes cyclically---that is, no index $a < b < N$ satisfies
$t_b.\pred = t_a.\pred$ in a way that allows positions $a, \ldots, b{-}1$
to form an independent closed sub-walk.
\end{definition}

\begin{lemma}[Enabling walks are simple --- no fragmentation under equivariance]
\label{lem:walk-simple}
Let $C$ be a simple cycle of length~$N$ in~$\Hstar(\Ls)$
at process~$P_r$ in a livelock on a ring of size~$K$.
Every enabling walk for~$C$ is a simple cycle of the
same length~$N$.
\end{lemma}

\begin{proof}
Let $W$ be an enabling walk for~$C$ at~$P_{r-1}$.
By Lemma~\ref{lem:equal-length}, $W$ is a closed walk
of length~$N$.
The shadow obligation chain of~$C$
(Definition~\ref{def:shadow-chain}) constrains~$W$ as
a single cyclic structure: the walk must satisfy all~$N$
shadow pairs in order and close.
We show that this structure cannot fragment into
independently closed components.

Suppose $W$ is not simple.
Then $W$ decomposes into simple cycles $S_1, \ldots, S_p$
with $|S_1| + \cdots + |S_p| = N$ and each $|S_i| < N$.
Each~$S_i$ independently satisfies closure: it is a valid
simple cycle in~$\Hstar(\Ls)$ with its own self-consistent
shadow obligations, satisfiable independently of~$C$.

Since $\Ls$ is a fixed point of~$\bPhi$, every simple
cycle in~$\Hstar(\Ls)$ admits enabling walks at every
process.
Therefore~$S_i$ propagates independently through
$P_{r-2}, P_{r-3}, \ldots$, sustained by equivariance:
each~$S_i$ carries its own shadow obligations and
produces a closed walk of length~$|S_i|$ at each
successive process.

The number of simple cycles of length~$|S_i|$
in~$\Hstar(\Ls)$ is at most $\binom{|\Ls|}{|S_i|}$.
By the pigeonhole principle, there exists $m \geq 1$ such
that after~$mK$ back-propagations, the image of~$S_i$
returns to~$P_r$ as a closed walk of length~$|S_i|$
whose transitions are drawn from~$\Ls$.

In the livelock execution, $P_r$ executes~$C$.
By Remark~\ref{rem:periodic}, the equivariant structure
at~$P_r$ is determined by~$C$.
The image of~$S_i$ at~$P_r$ is therefore a closed walk
of length $|S_i| < N$ that must be supported within~$C$.

But~$C$ is simple: it visits each transition exactly once,
and its shadow obligation chain is a single irreducible
cycle.
A closed sub-walk of length $|S_i| < N$ would constitute
an independent cyclic closure within~$C$---a proper
sub-sequence of the obligation chain that closes on its
own.
This contradicts the simplicity of~$C$, which admits no
proper closed sub-walk.

Therefore $W$ is simple.
\end{proof}

\begin{corollary}[Equivariant decomposition]
\label{cor:decomposition}
Any closed equivariant walk in a livelock decomposes into
simple cycles that individually satisfy closure and equivariance.
Moreover, each such simple cycle propagates bijectively
(up to cyclic rotation) to a simple cycle of the same length
at the predecessor process.
\end{corollary}

\begin{proof}
By Lemma~\ref{lem:decompose}, a closed walk contains a
simple cycle~$C$.
By Lemma~\ref{lem:walk-simple}, the enabling walk for~$C$
is a simple cycle of the same length.
The remaining transitions (if any) form a shorter closed walk,
to which the argument applies inductively.
\end{proof}

\begin{lemma}[Livelock reduces to simple-cycle circulation]
\label{lem:reduction}
If a livelock exists, then there exists a livelock supported
by a circulation of simple cycles, matched bijectively
(up to cyclic rotation) across processes.
All such cycles are contained in~$\Ls$.
In particular, $\Ls \neq \emptyset$.
\end{lemma}

\begin{proof}
By Remark~\ref{rem:periodic}, a livelock admits periodic
equivariant closed walks in~$\Hstar(\Ls)$.
By Lemma~\ref{lem:decompose}, each closed walk contains
a simple cycle.
By Lemma~\ref{lem:walk-simple}, the enabling walk for
this simple cycle is itself a simple cycle of the same length.
The transitions of these simple cycles form a set closed under
shadow back-propagation and pseudolivelock completion;
by Lemma~\ref{lem:prefixed}, this set is a pre-fixed point
of~$\bPhi$, hence contained in~$\Ls$.
\end{proof}

\subsection{Soundness:
\texorpdfstring{$\Ls \neq \emptyset$}{L* nonempty}
implies livelock}

\begin{lemma}[Simple cycles in $\Ls$ construct a livelock]
\label{lem:construction}
If $\Ls \neq \emptyset$, then there exists a simple cycle
in~$\Hstar(\Ls)$ that supports a livelock on any ring
whose size~$K$ satisfies the circulation law.
\end{lemma}

\begin{proof}
Since $\Ls = \nu\bPhi(T)$ is a nonempty fixed point of~$\bPhi$,
every transition in~$\Ls$ lies on a cycle
in~$\Hstar(\Ls)$.
By Lemma~\ref{lem:decompose}, this cycle contains a simple
cycle~$C$.
Assign to each process~$P_r$ the cycle~$C$, shifted by the
cyclic rotation determined by the enabling walk.
By Lemma~\ref{lem:walk-simple}, the enabling walk at each
predecessor is a simple cycle of the same length, hence
a valid execution cycle.
Each process repeatedly fires its transitions in order,
producing a periodic execution that never terminates.
The circulation law ensures that the cyclic rotations are
consistent across all~$K$ processes.
\end{proof}

\subsection{Main results}

\begin{theorem}[Semantic faithfulness]\label{thm:faithful}
For a self-disabling protocol with transition set~$T$:
\[
  \Ls \neq \emptyset
  \quad\Longleftrightarrow\quad
  \text{a livelock execution exists for some~$K \geq 2$.}
\]
\end{theorem}

\begin{proof}
$(\Leftarrow)$~Lemma~\ref{lem:reduction}.
$(\Rightarrow)$~Lemma~\ref{lem:construction}.
\end{proof}

\begin{theorem}[Bounded witness]\label{thm:bounded}
For a self-disabling protocol with transition set~$T$ on state
space~$\Zm$, it suffices to consider simple-cycle witnesses of
length at most~$|T|$: if a livelock exists for any ring size~$K$,
then it is witnessed by a circulation of simple cycles drawn
from~$\Ls$, each of length at most~$|T|$, independent of~$K$.
\end{theorem}

\begin{proof}
By Lemma~\ref{lem:reduction}, a livelock is supported by
simple cycles in~$\Hstar(\Ls)$.
A simple cycle visits each transition at most once, so its
length is at most~$|\Ls| \leq |T| \leq m^3$.
The number of ring circulations required for the witness
to close is at most~$\binom{|T|}{N}$ for a cycle of
length~$N$, by the pigeonhole argument in
Lemma~\ref{lem:walk-simple}.
Therefore the witness is bounded and its size depends only
on~$T$, not on~$K$.
\end{proof}

\begin{corollary}[Independent decidability]
\label{cor:decidability}
Livelock existence for self-disabling protocols on
unidirectional rings is decidable.
\end{corollary}

\begin{proof}
By Theorem~\ref{thm:bounded}, every livelock is witnessed
by simple cycles of length at most~$|T|$.
The set of such cycles is finite (at most
$\sum_{N=2}^{|T|} \binom{|T|}{N}$ candidates).
For each candidate, the enabling walk computation and
circulation law verification are finite procedures.
Livelock freedom holds if and only if no witness exists.
This decision procedure depends only on~$|T|$, not on~$K$.
\end{proof}

\begin{remark}[Relationship to the polynomial-time algorithm]
Corollary~\ref{cor:decidability} provides a decidability
proof that is independent of the fixed-point algorithm:
it relies solely on the bounded witness structure, not
on the computation of~$\Ls$.
The na\"ive procedure implied by the corollary---enumerating
simple cycles and checking the circulation law---has
complexity exponential in~$|T|$.
The $O(|T|^3)$ algorithm of the preceding section provides
a vastly more efficient decision procedure, but
Corollary~\ref{cor:decidability} demonstrates that
decidability follows from the structure of the problem
alone, without reference to any particular algorithm.
\end{remark}

\begin{remark}[Why self-disabling bounds the witness]
\label{rem:why}
The self-disabling property is the structural reason that
witness length is bounded.
In a self-disabling protocol, after firing
transition $(v, w, w')$, no transition $(v, w', {-})$
exists in~$T$; thus every continuation at the same process
requires a \emph{different} predecessor value.
This forces the predecessor values along any simple cycle
to form a sequence of distinct consecutive pairs (shadows),
of which there are at most $m(m{-}1)$.
Every simple cycle therefore has length at most~$m(m{-}1)$,
and the witness space is bounded.
Protocols that lack the self-disabling property can chain
transitions with identical predecessor values, admitting
unbounded local computations and hence unbounded witnesses.
\end{remark}

%% ═══════════════════════════════════════════════════════════════════════════
\section{Circulation Law}
\label{sec:circulation}
%% ═══════════════════════════════════════════════════════════════════════════

Farahat's original characterisation~\citep{farahat2012} formulates the
Circulation Law as $H^{|E|}\cap E^K\neq\emptyset$, an intersection
condition on the full binary relations $H$ and~$E$ over~$\Ls$.
The bounded witness theorem (Theorem~\ref{thm:bounded}) shows this
formulation is unnecessarily coarse: since every livelock reduces to
a circulation of simple cycles
(Corollary~\ref{cor:decomposition}), the Circulation Law should be
stated at the level of simple cycles, where it becomes explicit
arithmetic.

\begin{remark}[Equivariance on simple cycles]
\label{rem:equiv-simple}
By Lemma~\ref{lem:walk-simple}, equivariant propagation maps
each simple cycle of length~$n$ in $\Hstar(\Ls)$ bijectively
(up to cyclic rotation) to a simple cycle of the same length
at the successor process.
This bijective propagation of simple cycles is the operative content
of equivariance for livelock existence.
Farahat's equivariance condition $H\circ E = E\circ H$ on the full
binary relations $H$ and~$E$ is a sufficient algebraic encoding of
this bijective cycle propagation, but the bounded witness theorem
shows that the proper semantic level is simple cycles, where
propagation acts as a bijection with a cyclic shift.
Note that the propagation graph on simple cycles
(Definition~\ref{def:cl-simple}) is a general directed graph:
each cycle-to-cycle mapping is a bijection, but multiple source
cycles may share the same target, and non-deterministic shadow
witnesses may give a single source multiple outgoing edges.
\end{remark}

\begin{definition}[Circulation Law on simple cycles]
\label{def:cl-simple}
Let $\Ls\neq\emptyset$.
By Lemma~\ref{lem:walk-simple}, equivariant propagation maps
each simple cycle of length~$n$ in~$\Hstar(\Ls)$ bijectively
(up to cyclic rotation) to a simple cycle of the same length~$n$
in the successor process.
Each such mapping carries a cyclic shift~$s$.

The propagation defines a directed graph~$\mathcal{G}$ on
simple cycles: each node is a simple cycle in~$\Hstar(\Ls)$, and
each directed edge~$C\to C'$ carries a shift~$s$, meaning
that if $P_r$ executes~$C$, then $P_{r+1}$ executes~$C'$
with cyclic offset~$s$.
When the shadow map within~$\Ls$ has multiple witnesses,
a node may have multiple outgoing edges (non-deterministic
propagation); otherwise each node has a unique successor.

The \emph{Circulation Law} states: a livelock exists on a ring
of size~$K$ with~$e$ enabled processes if and only if
there exists a simple cycle~$C_0$ and a
\emph{closed walk} in~$\mathcal{G}$,
\[
  C_0 \to C_1 \to \cdots \to C_{qK-1} \to C_0,
\]
of length~$qK$ (for some $q\geq 1$), such that
\[
  e \;\equiv\; \sum_{i=0}^{qK-1} s_i \pmod{n},
  \qquad 0 < e \leq K,
\]
where:
\begin{itemize}
  \item $n$: the common length of every simple cycle in the walk
        (\textbf{constant}, by Lemma~\ref{lem:walk-simple}),
  \item $s_i$: the cyclic shift on the edge $C_i\to C_{i+1}$,
  \item $q \geq 1$: the number of complete ring circulations,
  \item $K \geq 2$: the ring size,
  \item $e$: the number of enabled processes (\textbf{derived}).
\end{itemize}
The walk may revisit cycles and need not be periodic.
The constraint $C_{qK} = C_0$ requires only that the walk
close after exactly~$qK$ steps; the intermediate cycles are
unconstrained beyond the graph structure of~$\mathcal{G}$.
\end{definition}

\begin{remark}[The Circulation Law is derived from equivariance]
\label{rem:cl-derived}
Definition~\ref{def:cl-simple} is not an independent condition on
livelock existence.
Given equivariance on the simple cycles of~$\Hstar(\Ls)$,
the graph~$\mathcal{G}$ is well-defined and non-empty
(since $\Ls\neq\emptyset$ implies at least one simple cycle).
Every node in~$\mathcal{G}$ has at least one outgoing edge,
so closed walks of arbitrary length exist.
For any starting cycle~$C_0$, the number of distinct simple
cycles of length~$n$ is finite (at most $\binom{|\Ls|}{n}$),
so closed walks returning to~$C_0$ exist for appropriate~$qK$.
The Circulation Law therefore determines \emph{which}
values of~$K$ and~$e$ are admissible, but never excludes
livelock existence entirely.
This is why the decision algorithm (Theorem~\ref{thm:main}) needs
only $\Ls\neq\emptyset$ and never checks the Circulation Law
explicitly: equivariance alone, captured by the fixed point~$\Ls$,
is necessary and sufficient for livelock existence.
\end{remark}

\begin{remark}[Computing admissible ring sizes]
For a given starting cycle~$C_0$ of length~$n$,
the admissible ring sizes are those~$K\geq 2$ for which
a closed walk of length~$qK$ exists in~$\mathcal{G}$
starting and ending at~$C_0$.
Since~$\mathcal{G}$ is a finite directed graph where every
node has at least one outgoing edge, closed walks of many
lengths exist, and the admissible~$(K, e)$ pairs are
computable from the graph structure.
\end{remark}

\begin{remark}[Two levels of analysis]
Simple cycles are the \emph{semantic} level at which livelock
structure is understood (Sections~\ref{sec:semantics}
and~\ref{sec:circulation}): equivariance acts bijectively,
the Circulation Law reduces to shift arithmetic, and witnesses
are bounded.
The full binary relations $H$ and~$E$ are the \emph{computational}
level at which the algorithm operates (Section~\ref{sec:correct}):
the fixed-point operator~$\bPhi$ works on transition sets via
relation-level shadow filtering, without enumerating simple cycles.
The bounded witness theorem bridges the two levels, guaranteeing
that the relation-level computation captures all cycle-level witnesses.
\end{remark}

%% ═══════════════════════════════════════════════════════════════════════════
\section{The $(1,1)$-Asymmetric Case}
%% ═══════════════════════════════════════════════════════════════════════════

By Definition~\ref{def:lq}, the $(1,1)$-asymmetric ring has $P_0$
running $T_0$ and $P_1,\ldots,P_{K-1}$ running $T_{\mathit{other}}$,
with valid ring sizes $K = 1+n$ for any $n\geq 1$.

A single application of Algorithm~1 to $T_{\mathit{other}}$ alone is
\emph{insufficient}: its fixed point $\Ls_{\mathit{other}}$ may contain
transitions whose shadow arcs are not supported by $T_0$.  Declaring
livelock based on $\Ls_{\mathit{other}}\neq\emptyset$ alone would be
unsound. Similarly, ring closure back to $P_0$ may shrink $L_0$, which
in turn would further constrain $\Ls_{\mathit{other}}$; this process
must iterate.

Algorithm~2 performs a \emph{joint fixed-point iteration} over the pair
$(L_0, L_{\mathit{other}})$:
\begin{enumerate}
\item Initialise $L_0 = \PL(T_0)$,
      $L_{\mathit{other}} = \PL(T_{\mathit{other}})$.
\item \emph{Constrain and re-stabilise $L_{\mathit{other}}$:} apply
      $\Filter(L_{\mathit{other}},\Shad(L_0))$ to remove transitions
      whose shadows $L_0$ cannot support, then re-run Algorithm~1
      warm-started from the result.
\item \emph{Ring closure:}
      $L_0' = \PL(\Filter(T_0,\Shad(L_{\mathit{other}})))$.
\item If either set empties, return \textsc{free}.
      If $L_0'=L_0$ (joint fixed point), return \textsc{livelock}.
      Otherwise set $L_0 \leftarrow L_0'$ and repeat.
\end{enumerate}

\paragraph{Termination.}
Both $L_0$ and $L_{\mathit{other}}$ are monotone decreasing across
iterations and bounded below by $\emptyset$.  By Knaster-Tarski on the
finite joint lattice $(2^{T_0}\!\times\!2^{T_{\mathit{other}}},\subseteq)$,
the iteration terminates in at most $|T_0|$ outer steps.

\paragraph{Complexity.}
With warm-starting, the total work for all inner fixed points across
all outer iterations is $O(|T_{\mathit{other}}|^3)$: each inner step
removes $\geq\!1$ transition from $L_{\mathit{other}}$, so the total
number of inner steps is at most $|T_{\mathit{other}}|$, each costing
$O(|T_{\mathit{other}}|^2)$.  Ring closures cost
$|T_0|\times O(|T_0|^2)=O(|T_0|^3)$.  Total:
$O(|T_{\mathit{other}}|^3+|T_0|^3)=O(|T|^3)$, independent of $K$.

\paragraph{Correctness.}
At the joint fixed point $(L_0^*, \Ls_{\mathit{other}})$, equivariance
holds at all three ring interfaces:
\begin{align*}
  H^*_{\mathit{other}}\circ E_{\mathit{cross}}
    &= E_{\mathit{cross}}\circ H_0^*
  & (P_{K-1}\to P_0) \\
  H_0^*\circ E_0
    &= E_0\circ H^*_{\mathit{other}}
  & (P_0 \to P_1) \\
  H^*_{\mathit{other}}\circ\Estar_{\mathit{other}}
    &= \Estar_{\mathit{other}}\circ H^*_{\mathit{other}}
  & (P_r \to P_{r+1})
\end{align*}
by the same Remark~\ref{rem:equiv} argument applied at each interface.
Lemma~\ref{lem:farahat} then gives a livelock.
Verified computationally for Dijkstra's token ring at $m\in\{3,4,5\}$
\citep{farahat2012}.

%% ═══════════════════════════════════════════════════════════════════════════
\section{Experimental Results}
%% ═══════════════════════════════════════════════════════════════════════════

Table~\ref{tab:results} shows results on standard self-stabilizing
protocols. Implementation:
\url{https://github.com/cosmoparadox/mathematical-tools}.

\begin{table}[h]
\centering
\caption{Results ($m=3$ unless noted). $|\Ls|=0$ means \textsc{free}.}
\label{tab:results}
\smallskip
\begin{tabular}{lcccc}
\toprule
Protocol & Type & $|\Ls|$ & Result & Min $K$ \\
\midrule
Dijkstra token ring    & $(1,1)$-asym & $3{+}3$  & \textsc{livelock} & $K\!\geq\!2$\\
$k$-Coloring (det)     & symmetric    & $m$      & \textsc{livelock} & $K\!\equiv\!1\pmod{m}$\\
$k$-Coloring (non-det) & symmetric    & $2m$     & \textsc{livelock} & $K\!\geq\!3$\\
Agreement (maximal)    & symmetric    & $m(m{-}1)$ & \textsc{livelock} & $K\!\geq\!2$\\
Sum-Not-2 (det)        & symmetric    & $0$      & \textsc{free}     & ---\\
Sum-Not-2 (non-det)    & symmetric    & $2m$     & \textsc{livelock} & $K\!\geq\!2$\\
\bottomrule
\end{tabular}
\end{table}

Note: The livelock detected for Dijkstra's token-ring is due to the fact that the specification is for a constant space of size $m$ per process. The ring will have a single token circulation for ring sizes smaller than $1 + m$. Multiple token circulations show for $K \geq 1 + m$. Our implementation of the algorithm detects that pattern through the Circulation Law.
%% ═══════════════════════════════════════════════════════════════════════════
\section{Related Work}
%% ═══════════════════════════════════════════════════════════════════════════

\citet{klinkhamer2013,klinkhamer2019} establish $\Pi^0_1$-completeness
for parameterized livelock freedom and provide a semi-algorithm for
livelock \emph{detection}. The authors characterize livelock existence via periodic tiling patterns capturing local propagation constraints. In contrast, the present work establishes the boundedness of the livelock witnesses establishing a baseline for decidability. In addition, we enforce closure of recurrent supports under pseudolivelocks and propagation via a fixed-point construction. This contributes a polynomial decision procedure for livelock \emph{freedom}, grounded in Farahat's
algebraic framework~\citep{farahat2012}: the greatest fixed point $\Ls$
of $\bPhi$ decides freedom for all $K$ in $O(|T|^3)$ time. This suggests that tiling-based formulations may admit compound cyclic patterns that satisfy local consistency but do not correspond to fully realizable recurrent supports in livelocks of self-disabling rings.

\citet{suzuki1988} proves undecidability of general temporal property
verification for rings of finite-state machines. Our result does not
contradict this: livelock detection is a specific structural property
admitting the $O(|T|^3)$ procedure of Theorem~\ref{thm:main} for the stricter execution semantics of self-disabling processes.

\citet{emerson1995} establish cutoff results for correctness properties
on token rings. Our approach is complementary: we eliminate the
dependence on $K$ entirely by computing $\Ls$ from $T$.

%% ═══════════════════════════════════════════════════════════════════════════
\section{Conclusion}
%% ═══════════════════════════════════════════════════════════════════════════

Livelock detection for parameterized self-disabling unidirectional rings
is decidable in $O(|T|^3)$ time independent of $K$. The decision is
$\Ls\neq\emptyset$. Livelock freedom is certified by the maximality of
$\Ls$: Invariant~\ref{inv:max} guarantees $\Hi{i}$ is maximal at every
iterate; when $\Ls=\emptyset$, no process can sustain a pseudolivelock
for any $K$.
Moreover, the bounded witness theorem (Theorem~\ref{thm:bounded})
provides an independent proof of decidability by showing that every
livelock witness---regardless of ring size---reduces to simple cycles
of length at most~$|T|$.
The Circulation Law, which determines admissible ring sizes, is a
derived consequence of equivariance on simple cycles
(Remark~\ref{rem:cl-derived}): valid ring sizes always exist given
a nonempty~$\Ls$, so the algorithm never needs to check the
Circulation Law explicitly.
Extension to the general $(l,q)$-asymmetric case ($K=l+nq$) is left to
a companion paper. Code and algebraic foundation:
\url{https://github.com/cosmoparadox/mathematical-tools}.

%% ── references ────────────────────────────────────────────────────────────
\bibliographystyle{plainnat}
\begin{thebibliography}{9}

\bibitem[Dijkstra(1974)]{dijkstra1974}
E.~W. Dijkstra.
\newblock Self-stabilizing systems in spite of distributed control.
\newblock \emph{Communications of the ACM}, 17\penalty0 (11):\penalty0
  643--644, 1974.

\bibitem[Emerson and Namjoshi(1995)]{emerson1995}
E.~A. Emerson and K.~S. Namjoshi.
\newblock Reasoning about rings.
\newblock In \emph{POPL}, pages 85--94, 1995.

\bibitem[Farahat(2012)]{farahat2012}
A.~Farahat.
\newblock \emph{Automated Design of Self-Stabilization}.
\newblock PhD thesis, Michigan Technological University, 2012.
\newblock Chapter~6: Exact Algebraic Characterization of Global Convergence.
  Code and chapter at
  \url{https://github.com/cosmoparadox/mathematical-tools}.

\bibitem[Klinkhamer and Ebnenasir(2013)]{klinkhamer2013}
A.~P. Klinkhamer and A.~Ebnenasir.
\newblock Verifying livelock freedom on parameterized rings and chains.
\newblock In \emph{SSS}, pages 163--177, 2013.

\bibitem[Klinkhamer and Ebnenasir(2019)]{klinkhamer2019}
A.~P. Klinkhamer and A.~Ebnenasir.
\newblock Verifying livelock freedom on parameterized rings and chains.
\newblock \emph{ACM Transactions on Computational Logic}, 20\penalty0
  (3):\penalty0 1--49, 2019.

\bibitem[Suzuki(1988)]{suzuki1988}
I.~Suzuki.
\newblock Proving properties of a ring of finite-state machines.
\newblock \emph{Information Processing Letters}, 28\penalty0 (4):\penalty0
  213--214, 1988.

\bibitem[Tarski(1955)]{tarski1955}
A.~Tarski.
\newblock A lattice-theoretical fixpoint theorem and its applications.
\newblock \emph{Pacific Journal of Mathematics}, 5\penalty0 (2):\penalty0
  285--309, 1955.

\end{thebibliography}

\end{document}
